\section{Retention of every prescribed maximum simplex}
\label{sec:retention}
This section reproduces the retained retention estimates with the quantifier
on the original simplex explicit. The arguments use only containment and
near-vertex points, independently of how those points were obtained.

\begin{lemma}[Substochastic determinant matching]\label{lem:matching}
Let $U$ be a nonnegative $l\times l$ column-substochastic matrix, where
$l\ge1$. If $|\det U|\ge1-\eta$ and $0\le\eta<1/3$, then a permutation
of its columns makes every diagonal entry at least $1-\eta$.
\end{lemma}
\begin{proof}
Represent each column as the expectation of a random standard basis
vector, assigning any missing mass to the zero vector. Independent
selection represents the determinant as the expectation of a variable in
$\{-1,0,1\}$. If every column but column $j$ is sampled, each resulting
cofactor vector either vanishes or selects exactly one entry of column $j$
with a sign. Hence $|\det U|\le\max_iU_{ij}$. Choose an entry at least
$1-\eta$ in each column. If two chosen entries share a row, the independent
selections collide with probability at least $(1-\eta)^2$, so
\[
 |\det U|\le1-(1-\eta)^2\le2\eta<1-\eta,
\]
a contradiction. The selected rows are distinct.
\end{proof}

\begin{lemma}[Inverse bootstrap]\label{lem:bootstrap}
Let $K\subset P$, where $P$ is a $d$-simplex and $S$ is a prescribed
maximum-volume simplex in $K$. Suppose $q_i\in K$ have $P$-coordinates
$Q_{ii}\ge1-\rho$. Suppose that after relabeling the vertices of this same
$S$, their $P$-coordinate matrix $W$ satisfies $W_{ii}\ge1-\eta$, where
$0\le\eta\le1/8$ and $0\le\rho\le1/16$. Then
\begin{equation}\label{eq:bootstrap}
 E(K,S)\le8(d+1)\rho.
\end{equation}
\end{lemma}
\begin{proof}
The matrix $W$ is nonnegative column stochastic and invertible. Put
$B=W^{-1}$, $D=B-I$, $\beta=\max_{ij}|D_{ij}|$, and
$r_j=1-W_{jj}$. From $DW=I-W$ and column stochasticity,
\[
 D_{ij}=(\delta_{ij}-W_{ij})+
           \sum_{l\ne j}(D_{ij}-D_{il})W_{lj}.
\]
Consequently $\beta\le\eta+2\eta\beta$, whence
\begin{equation}\label{eq:firstinverse}
 \beta\le\frac\eta{1-2\eta}\le\frac16.
\end{equation}
The $ii$ entry of $BW=I$ gives
\[
 B_{ii}-1\ge\frac{(1-\beta)r_i}{1-r_i}\ge r_i/2.
\]
On the other hand, $q_i\in K$ and maximality of the same $S$ imply
$(BQ)_{ii}\le1$ by replacement. Since $Q_{ii}\ge1-\rho$,
$B_{ii}>0$, and $B_{ij}\ge-\beta$ for $i\ne j$,
\[
 B_{ii}-1\le\rho(B_{ii}+\beta)\le2\rho.
\]
Thus $r_i\le4\rho$ for every $i$. Reapply the inverse identity with this
improved column loss to obtain
\[
 \beta\le\frac{4\rho}{1-8\rho}\le8\rho.
\]
Every $S$-coordinate of each $P$-vertex is an entry of $B$ and is at least
$-8\rho$. Convex averaging gives the same lower bound throughout $P$.
Equation~\eqref{eq:dilationcoords} now yields \eqref{eq:bootstrap} about
$S$'s own centroid.
\end{proof}

\begin{proposition}[Two all-scale retention bounds]\label{prop:retention}
Let $d\ge1$ be an integer, $n=d+1$, and let $K\subset P$, where $K$ is a
full-dimensional compact convex body and $P$ a $d$-simplex. Suppose every
$P$-coordinate attains at least $1-\rho$ on $K$, where $\rho\ge0$.
For every prescribed maximum-volume inscribed simplex $S$,
\begin{equation}\label{eq:unrestrictedretention}
 E(K,S)\le16n^2\rho.
\end{equation}
If, additionally, $K$ is a polytope with at most $n+r$ vertices for an
integer $r\ge1$, then
\begin{equation}\label{eq:boundedretention}
 E(K,S)\le64nr\rho.
\end{equation}
The local estimate $E(K,S)\le8n\rho$ holds in the first case for
$\rho\le1/(16n)$ and in the second case for $\rho\le1/(64r)$.
\end{proposition}
\begin{proof}
The universal estimate $E(K,S)\le n$ follows from maximum-volume
replacement and \eqref{eq:dilationcoords}, without any normalization.
It handles the complementary large-$\rho$ branches of both bounds.

For the unrestricted local branch, choose $q_i\in K$ with $P$-coordinate
at least $1-\rho$ and let $Q$ be their column coordinate matrix. Independent
column selection gives
\[
 \det Q\ge2\prod_iQ_{ii}-1\ge1-2n\rho.
\]
Indeed the identity selection has probability $\prod_iQ_{ii}$ and every
other signed selection is at least $-1$. For $\rho\le1/(16n)$ the right
side is positive, so these points form a comparison simplex. Maximality of
$S$ gives $|\det W|\ge\det Q\ge1-\eta$ with $\eta=2n\rho\le1/8$.
Lemma~\ref{lem:matching} relabels only the vertices of this same $S$ so
that $W_{ii}\ge1-\eta$. Lemma~\ref{lem:bootstrap} proves the local bound.
Together with $E\le n$ for $\rho>1/(16n)$, this proves
\eqref{eq:unrestrictedretention}.

For the bounded-vertex branch assume $\rho\le1/(64r)$. Choose actual
vertices $q_i$ maximizing the $P$-coordinates. They are distinct because
$\rho<1/2$. Their nonnegative column-stochastic matrix satisfies
\begin{equation}\label{eq:Qinverse}
 \norm{Q-I}_1\le2\rho,
 \qquad \norm{Q^{-1}}_1\le(1-2\rho)^{-1}\le1+4\rho.
\end{equation}
Here $\norm{\cdot}_1$ is the induced column-sum matrix norm and
$\rho\le1/4$. Thus $T=\conv(q_i)$ is a simplex in $K$.

First let $S_v$ be any maximum simplex with vertices among those of $K$.
Such a maximum exists by multiaffinity of the determinant. It omits at most
$r$ of the $q_i$ because there are at most $r$ other vertices. Let $R$ be
the omitted index set, of size $l\le\min(r,n)$. If $l=0$, then $S_v=T$.
Otherwise label its retained $q_i$ by their own indices and put the remaining
vertices in columns $R$. If $W$ is its coordinate matrix and $C_0=Q^{-1}W$,
then $(C_0)_i=e_i$ for $i\notin R$, so
\[
 |\det (C_0)_{R,R}|=|\det C_0|=\frac{|S_v|}{|T|}\ge1.
\]
The restricted matrix $U=W_{R,R}$ is column substochastic. Each column of
$(C_0)_{R,R}$ has $\ell^1$ norm at most $1+4\rho$; by
$W-C_0=(Q-I)C_0$ its distance from the corresponding $U$-column is at most
$2\rho/(1-2\rho)\le4\rho$. Telescoping the determinant gives
\[
 \bigl|\det U-\det(C_0)_{R,R}\bigr|
 \le4l\rho(1+4\rho)^{l-1}\le8l\rho.
\]
For the last inequality, $4l\rho\le1/16$ and
$(1+4\rho)^{l-1}\le(1-4l\rho)^{-1}<2$, by the binomial expansion.
Hence $|\det U|\ge1-\eta$ with $\eta=8r\rho\le1/8$.
Lemma~\ref{lem:matching} matches the omitted indices, while the retained
$q_i$ already have diagonal coordinates at least $1-\rho\ge1-\eta$.
Every maximum vertex simplex therefore has exactly one vertex in each of
the disjoint convex caps
\begin{equation}\label{eq:disjointcaps}
 \mathcal C_i=\{x\in K:\lambda_i^P(x)\ge1-\eta\}.
\end{equation}
This conclusion also covers $l=0$.

Now return to the prescribed maximum $S$, whose vertices need not be
vertices of $K$. Express each vertex of $S$ as a convex combination of
$K$-vertices and sample the combinations independently. Its determinant is
the expectation of the sampled determinants. Every sampled absolute
volume is at most $|S|$. Equality in the expectation and triangle
inequality forces every positive-probability sampled tuple to be a maximum
vertex simplex, with one point in each cap \eqref{eq:disjointcaps}.
For each $i$, some sampled column must lie in $\mathcal C_i$ with probability
one; otherwise independence would give positive probability that all
columns avoid it. Distinct caps require distinct such columns. There are
$n$ columns and $n$ caps, so after relabeling, each actual vertex of this
same $S$ lies in its corresponding cap by convexity.

Thus $W_{ii}\ge1-\eta$ for the prescribed $S$, and
Lemma~\ref{lem:bootstrap} proves $E(K,S)\le8n\rho$ locally.
For $\rho>1/(64r)$ the universal estimate gives $E(K,S)<64nr\rho$.
No change of simplex carrier or centroid has occurred.
\end{proof}
